

\[ T_R = T_R(m_1, m_2, r) \]

\noindent Identificando as variáveis e dimensões:
\begin{align*}
T_R \text{ (tempo):} \quad &[T] \\
m_1 \text{ (massa do corpo 1):} \quad &[M] \\
m_2 \text{ (massa do corpo 2):} \quad &[M] \\
r \text{ (distância entre eles):} \quad &[L]
\end{align*}

Temos $n = 4$ variáveis e $m = 3$ dimensões. Logo, $n - m = 4 - 3 = 1$ grupo adimensional.
\begin{align*}
\pi &= (m_1)^a \cdot (m_2)^b \cdot (r)^c \cdot (T_R)^d \\
\pi &= (M)^a \cdot (M)^b \cdot (L)^c \cdot (T)^d
\end{align*}

Analisando os expoentes:
\[
\begin{cases}
\text{Para M: } a + b = 0 \implies a = 1, b = -1 \implies \frac{m_1}{m_2} \quad \text{(ok!)} \\
\text{Para L: } c = 0 \quad \text{(o comprimento some)} \\
\text{Para T: } d = 0 \quad \text{(o tempo também some)}
\end{cases}
\]

Ao tentar formar o grupo Pi, percebe-se que as dimensões $L$ (comprimento, referente a $r$) e $T$ (tempo, referente a $T_R$) estão isoladas. Aparentemente a equação (2.14) está incompleta. É necessária uma ferramenta (constante) que relacione tempo e espaço e converta massa e distância em tempo. 

A constante gravitacional $G$, que possui as dimensões $[L^3 M^{-1} T^{-2}]$, pode ser usada para que exista uma relação dimensional homogênea entre massa, tempo e distância no problema de dois corpos.

\vspace{0.5cm}
\noindent\textbf{Refazendo com a adição de $G$:}
\[ T_R = T_R(m_1, m_2, r, G) \]
Adiciona-se $G$ (constante gravitacional): $[L^3 M^{-1} T^{-2}]$. \\
Agora temos $n = 5$ variáveis e $m = 3$ dimensões. \\
Grupos $\pi = n - m = 5 - 3 = 2$.

O primeiro grupo já sabemos que é a razão das massas:
\[ \pi_1 = \frac{m_1}{m_2} \]

Para o segundo grupo, fixamos o expoente de $T_R$ como 1:
\begin{align*}
\pi_2 &= G^a \cdot m_1^b \cdot r^c \cdot T_R^1 \\
M^0 L^0 T^0 &= (L^3 M^{-1} T^{-2})^a \cdot (M)^b \cdot (L)^c \cdot (T)^1
\end{align*}

Montando o sistema:
\[
\begin{cases}
\text{T: } -2a + 1 = 0 \implies a = 1/2 \\
\text{M: } -a + b = 0 \implies b = 1/2 \\
\text{L: } 3a + c = 0 \implies 3(1/2) + c = 0 \implies c = -3/2
\end{cases}
\]

Substituindo na expressão de $\pi_2$:
\begin{align*}
\pi_2 &= G^{1/2} \cdot m_1^{1/2} \cdot r^{-3/2} \cdot T_R \\
\pi_2 &= T_R \sqrt{\frac{G \cdot m_1}{r^3}}
\end{align*}

Elevando ao quadrado para simplificar:
\[ \pi_2^2 = \frac{T_R^2 \cdot G \cdot m_1}{r^3} = \text{constante} \]
\[ T_R^2 \propto r^3 \]

\textbf{Conclusão:} A Terceira Lei de Kepler, ao incluir a constante $G$, previu que o quadrado do tempo de revolução é proporcional ao cubo da distância.